% TOP_PROBLEM: {"id":30007568,"problem_number":"LOCAL-30007568","title":"Strategic queue stability","category_id":21,"category":{"id":21,"name":"economics","display_name":"Economics","description":"Open economic theory statements, published research agendas, and proposed scoped specifications; question provenance and historical priority are qualified per record.","slug":"economics","order_index":21,"created_at":"2026-10-08T00:00:00Z"},"status":"research_question","external_url":null,"legacy_ids":[],"published":false,"statement":"Determine the service-capacity augmentation needed for strong stability of strategic learning queues under specified compatibility graphs, buffers, and admission policies.","economics":{"original_id":57,"field":"Game theory and strategic interaction","kind":"theoretical","formalization_basis":"adapted_research_question","source_status":"research_agenda","priority_status":"earliest_found","priority_note":"This is the earliest source in which the relevant agenda was verified during this audit; absolute priority has not been established.","earliest_verified_reference":{"authors":["Éva Tardos"],"title":"Learning and the Price of Anarchy in Games","year":2026,"url":"https://epubs.siam.org/doi/full/10.1137/25M1808064","doi":"10.1137/25M1808064","locator":"Learning outcomes in dynamic games; Packet routing as a dynamic game; Theorem8","evidence_summary":"The lecture identifies learning with carry-over states as an active open agenda while presenting established capacity-augmentation results for queues."},"current_reference":{"authors":["Éva Tardos"],"title":"Learning and the Price of Anarchy in Games","year":2026,"url":"https://epubs.siam.org/doi/full/10.1137/25M1808064","doi":"10.1137/25M1808064","locator":"Learning outcomes in dynamic games; Packet routing as a dynamic game; Theorem8","evidence_summary":"The lecture identifies learning with carry-over states as an active open agenda while presenting established capacity-augmentation results for queues."},"formalization_note":"This is an adapted capacity frontier, not a claim that strategic queues or all network extensions lack stability theorems."},"statusQualification":"Published question or research agenda verified at the cited locator. The mathematical specification is adapted for this catalogue and includes additional assumptions; source publication does not establish first-ever priority or certify this exact model remains unsolved. This is an adapted capacity frontier, not a claim that strategic queues or all network extensions lack stability theorems.","statusReviewedAt":"2026-10-08"}
% FIRST_POSTED: 2026-10-08
% ENTRY_KIND: statement_only
% !TeX program = lualatex
\documentclass[11pt]{article}
\usepackage{fontspec}
\usepackage[a4paper,margin=25mm]{geometry}
\usepackage{amsmath,amssymb,mathtools}
\usepackage{xurl}
\usepackage[hidelinks]{hyperref}
\newcommand{\sref}[2]{\hyperlink{source:#1:#2}{[S#2]}}
\setlength{\emergencystretch}{3em}
\begin{document}
% problemId:30007568
\section*{Strategic queue stability}\label{30007568}
\subsection{Definitions and mathematical statement}
Let \(n\) source queues receive independent Bernoulli arrivals with rates \(\lambda_i\in(0,1)\). Each nonempty queue selects one compatible server per round in a bipartite compatibility graph \(E\). Server \(j\) processes an admitted packet with independent probability \(\mu_j\in(0,1]\) and has buffer capacity \(b_j\); a specified admission policy resolves collisions and returned packets remain at their source. Let \(Q_i(t)\) count source packets and \(H_j(t)\) buffered packets. Call the system strongly stable if
\[
 \sup_{t\geq0}E\left[\sum_iQ_i(t)+\sum_jH_j(t)\right]<\infty.
\]
Fix a high-probability, finite-window regret requirement for each queue's routing learner and specify its reward and feedback. Define \(c^*(E,b,\mathcal L)\) as the least service-capacity augmentation factor that guarantees stability for every centrally feasible arrival vector and every learner profile in \(\mathcal L\), under the stated admission policy. Determine this factor, useful instance-dependent bounds, and necessary restrictions on arrivals or buffer policies. Central feasibility must use the same compatibility graph and operational constraints, rather than total capacity alone. Compare matching constructive guarantees with destabilizing learner or collision examples. Factor-two and factor-three guarantees already proved for particular models are maintained; the target is the remaining architecture-dependent characterization, and each claimed numerical gap requires a fresh check against the literature.

\par\textbf{Formulation and scope.} This is an adapted capacity frontier, not a claim that strategic queues or all network extensions lack stability theorems.

\par\textbf{Open-question provenance.} An explicit question or research agenda was verified at the source locator. This does not by itself certify that every later version remains unresolved \sref{30007568}{1}.

\par\textbf{Historical priority.} This is the earliest source in which the relevant agenda was verified during this audit; absolute priority has not been established.

\subsection{Short English statement}
Determine the service-capacity augmentation needed for strong stability of strategic learning queues under specified compatibility graphs, buffers, and admission policies.

\subsection{Sources}
\begin{itemize}\raggedright
\item[\textnormal{[S1]}]\hypertarget{source:30007568:1}{} Éva Tardos. 2026. Learning and the Price of Anarchy in Games. Learning outcomes in dynamic games; Packet routing as a dynamic game; Theorem8.
\url{https://epubs.siam.org/doi/full/10.1137/25M1808064}.
DOI: 10.1137/25M1808064.
{\small\textit{Earliest verified question/agenda reference; historical priority qualified below. The lecture identifies learning with carry-over states as an active open agenda while presenting established capacity-augmentation results for queues.}}
\end{itemize}
\end{document}
